% GENERATED FILE. Edit canonical lessons, then run npm run generate:books.
% canonical-source-hash: fnv1a64:1da00fa0429c6546
% canonical-lessons: AR-C142-darraja, AR-C142-safina, AR-C142-musafir, AR-C142-masna, AR-C142-saha

\chapter{Bicycle, Ship, Traveller, Factory, Square}
\label{ch:ar-bicycle-ship-traveller-factory-square}
\hlchaptermodality{142}

\begin{chapteropening}
\textbf{By the end of this chapter:} \emph{I can say a bicycle, a ship, a traveller, a factory and a square.}

Complete the last lesson of chapter 142: I can say a bicycle, a ship, a traveller, a factory and a square.
\end{chapteropening}

\section[\texorpdfstring{\ar{دراجة} (darrāja)}{دراجة (darrāja)}]{\ar{دراجة} (darrāja) — a bicycle}

\label{lesson:AR-C142-darraja}



\begin{quote}
Before the new one: say the Arabic for skin, then the Arabic for a bone.
\end{quote}

\subsection*{You'll want to know: \ar{دراجة}}
\textbf{\ar{دراجة}} — \emph{darrāja} — \textquotedblleft{}a bicycle\textquotedblright{}.

\begin{cousinweb}
\textbf{\ar{دراجة}} (\emph{darrāja}) is built on the native Arabic root \emph{d-r-j}, \textquotedblleft{}to go step by step\textquotedblright{} (Lane). Lane's \emph{darrāja} is a wheeled frame on which a child or an old man steps along; the bicycle took over the word.
\end{cousinweb}

\begin{grammarlens}[title={What you've built}]
One more word for this chapter.
\end{grammarlens}

\subsection*{Your turn}
\begin{itemize}
\raggedright
  \item \emph{Say it:} \emph{darrāja}
  \item \emph{Say it:} \emph{darrāja}, once more
  \item \emph{Say it:} say \emph{darrāja}
  \item \emph{Recall:} say the Arabic for skin, then the Arabic for a bone, then say \emph{darrāja} again
\end{itemize}

\begin{tcolorbox}[breakable,colback=teal!4,colframe=teal!35!black,title={Before you move on}]
What does \ar{دراجة} mean? (\textquotedblleft{}A bicycle\textquotedblright{}.) Say it once more.
\end{tcolorbox}
\section[\texorpdfstring{\ar{سفينة} (safīna)}{سفينة (safīna)}]{\ar{سفينة} (safīna) — a ship}

\label{lesson:AR-C142-safina}



\begin{quote}
Before the new one: say the Arabic for a bone, then the Arabic for a bicycle.

Colours: what does the feminine do to \emph{aswad}, \emph{abyaḍ}, \emph{aḥmar} and \emph{azraq}? (Drops the opening \textbf{\ar{أ}} and ends in \emph{-āʾ}.)
\end{quote}

\subsection*{You'll want to know: \ar{سفينة}}
\textbf{\ar{سفينة}} — \emph{safīna} — \textquotedblleft{}a ship\textquotedblright{}.

\begin{cousinweb}
\textbf{\ar{سفينة}} (\emph{safīna}) is built on the native Arabic root \emph{s-f-n}, \textquotedblleft{}a ship; to pare\textquotedblright{} (Lane). Lane gives \emph{safīna} as a ship.
\end{cousinweb}

\begin{grammarlens}[title={What you've built}]
One more word for this chapter.
\end{grammarlens}

\subsection*{Your turn}
\begin{itemize}
\raggedright
  \item \emph{Say it:} \emph{safīna}
  \item \emph{Say it:} \emph{safīna}, once more
  \item \emph{Say it:} say \emph{safīna}
  \item \emph{Recall:} say the Arabic for a bone, then the Arabic for a bicycle, then say \emph{safīna} again
\end{itemize}

\begin{tcolorbox}[breakable,colback=teal!4,colframe=teal!35!black,title={Before you move on}]
What does \ar{سفينة} mean? (\textquotedblleft{}A ship\textquotedblright{}.) Say it once more.
\end{tcolorbox}
\section[\texorpdfstring{\ar{مسافر} (musāfir)}{مسافر (musāfir)}]{\ar{مسافر} (musāfir) — a traveller}

\label{lesson:AR-C142-musafir}



\begin{quote}
Before the new one: say the Arabic for a bicycle, then the Arabic for a ship.
\end{quote}

\subsection*{You'll want to know: \ar{مسافر}}
\textbf{\ar{مسافر}} — \emph{musāfir} — \textquotedblleft{}a traveller\textquotedblright{}.

\begin{cousinweb}
\textbf{\ar{مسافر}} (\emph{musāfir}) is built on the native Arabic root \emph{s-f-r}, \textquotedblleft{}to travel; to uncover\textquotedblright{} (Lane). \emph{Musāfir} is the Form III participle, \textquotedblleft{}one who travels\textquotedblright{}.
\end{cousinweb}

\begin{grammarlens}[title={What you've built}]
One more word for this chapter.
\end{grammarlens}

\subsection*{Your turn}
\begin{itemize}
\raggedright
  \item \emph{Say it:} \emph{musāfir}
  \item \emph{Say it:} \emph{musāfir}, once more
  \item \emph{Say it:} say \emph{musāfir}
  \item \emph{Recall:} say the Arabic for a bicycle, then the Arabic for a ship, then say \emph{musāfir} again
\end{itemize}

\begin{tcolorbox}[breakable,colback=teal!4,colframe=teal!35!black,title={Before you move on}]
What does \ar{مسافر} mean? (\textquotedblleft{}A traveller\textquotedblright{}.) Say it once more.
\end{tcolorbox}
\section[\texorpdfstring{\ar{مصنع} (maṣnaʿ)}{مصنع (maṣnaʿ)}]{\ar{مصنع} (maṣnaʿ) — a factory}

\label{lesson:AR-C142-masna}



\begin{quote}
Before the new one: say the Arabic for a ship, then the Arabic for a traveller.
\end{quote}

\subsection*{You'll want to know: \ar{مصنع}}
\textbf{\ar{مصنع}} — \emph{maṣnaʿ} — \textquotedblleft{}a factory\textquotedblright{}.

\begin{cousinweb}
\textbf{\ar{مصنع}} (\emph{maṣnaʿ}) is built on the native Arabic root \emph{ṣ-n-ʿ}, \textquotedblleft{}to make, to fashion\textquotedblright{} (Lane). \emph{Maṣnaʿ} is the \emph{ma-} place noun, \textquotedblleft{}a place of making\textquotedblright{}.
\end{cousinweb}

\begin{grammarlens}[title={What you've built}]
One more word for this chapter.
\end{grammarlens}

\subsection*{Your turn}
\begin{itemize}
\raggedright
  \item \emph{Say it:} \emph{maṣnaʿ}
  \item \emph{Say it:} \emph{maṣnaʿ}, once more
  \item \emph{Say it:} say \emph{maṣnaʿ}
  \item \emph{Recall:} say the Arabic for a ship, then the Arabic for a traveller, then say \emph{maṣnaʿ} again
\end{itemize}

\begin{tcolorbox}[breakable,colback=teal!4,colframe=teal!35!black,title={Before you move on}]
What does \ar{مصنع} mean? (\textquotedblleft{}A factory\textquotedblright{}.) Say it once more.
\end{tcolorbox}
\section[\texorpdfstring{\ar{ساحة} (sāḥa)}{ساحة (sāḥa)}]{\ar{ساحة} (sāḥa) — a square}

\label{lesson:AR-C142-saha}



\begin{quote}
Before the new one: say the Arabic for a traveller, then the Arabic for a factory.
\end{quote}

\subsection*{You'll want to know: \ar{ساحة}}
\textbf{\ar{ساحة}} — \emph{sāḥa} — \textquotedblleft{}a square\textquotedblright{}.

\begin{cousinweb}
\textbf{\ar{ساحة}} (\emph{sāḥa}) is built on the native Arabic root \emph{s-w-ḥ}, \textquotedblleft{}an open court\textquotedblright{} (Lane). Lane's \emph{sāḥa} is the open space of a house.
\end{cousinweb}

\begin{grammarlens}[title={What you've built}]
One more word for this chapter.
\end{grammarlens}

\subsection*{Your turn}
\begin{itemize}
\raggedright
  \item \emph{Say it:} \emph{sāḥa}
  \item \emph{Say it:} \emph{sāḥa}, once more
  \item \emph{Say it:} say \emph{sāḥa}
  \item \emph{Recall:} say the Arabic for a traveller, then the Arabic for a factory, then say \emph{sāḥa} again
\end{itemize}

\begin{tcolorbox}[breakable,colback=teal!4,colframe=teal!35!black,title={Before you move on}]
What does \ar{ساحة} mean? (\textquotedblleft{}A square\textquotedblright{}.) Say it once more.
\end{tcolorbox}
